% Part: normal-modal-logic
% Chapter: filtrations
% Section: preliminaries

\documentclass[../../../include/open-logic-section]{subfiles}

\begin{document}

\olfileid{nml}{fil}{pre}

\olsection{ప్రాథమిక విషయాలు}

వడపోతల ద్వారా మన మోడల్ తర్క వ్యవస్థలకు
\emph{పరిమిత నమూనా ధర్మం} ఉందని చూపి వాటి
నిర్ణేయతను స్థాపించవచ్చు. అంటే, ఒక నమూనాలో సత్యమైన
(అసత్యమైన) ఏ \tetoken{సూత్రమైనా}{!!{formula}} ఒక
\emph{పరిమిత} నమూనాలో కూడా సత్యం (అసత్యం) అవుతుంది.
ఉపసూత్రాల పరంగా సంవృతమైన
\tetoken{సూత్రాల}{!!{formula}s} సమితుల సాపేక్షంగా
వడపోతలను నిర్వచిస్తాం.

\begin{defn}\ollabel{defn:modallyclosed}
  \tetoken{సూత్రాల}{!!{formula}s} సమితి $\Gamma$లోని
  \tetoken{ఒక సూత్రానికి}{!!a{formula}} చెందిన ప్రతి
  ఉపసూత్రం కూడా $\Gamma$లో ఉంటే, $\Gamma$ను
  \emph{ఉపసూత్రాల పరంగా సంవృతం} అంటాం. అదనంగా అది
  ఉపసూత్రాల పరంగా సంవృతమై, $!A \in \Gamma$ అయినప్పుడల్లా
  $\Box!A, \Diamond!A \in \Gamma$ కూడా అయితే,
  $\Gamma$ను \emph{మోడల్ సంయోజకాల పరంగా సంవృతం} అంటాం.
\end{defn}

ఉదాహరణకు, \tetoken{ఒక సూత్రం}{!!a{formula}}~$!A$
ఇచ్చినప్పుడు, దాని అన్ని
\tetoken{ఉపసూత్రాల}{!!{formula}s} సమితి
\tetoken{ఉపసూత్రాల}{!!{formula}s} పరంగా సంవృతం.
ఆ \tetoken{ఉపసూత్రాల}{!!{formula}s} సమితి ద్వారా
ఒక నమూనా వడపోతను నిర్వచిస్తే, మనకు కావలసిన ధర్మం
దానికి ఉంటుంది: మూల నమూనాలో $!A$ సత్యం (అసత్యం)
అయితే, అప్పుడు మాత్రమే వడపోతలోనూ $!A$ సత్యం
(అసత్యం) అవుతుంది.

$\Gamma$ ద్వారా $\mModel{M}$కు చేసే వడపోతలోని లోకాల
సమితిని, కింది తుల్యతా సంబంధానికి చెందిన అన్ని
తుల్యతా వర్గాల సమితిగా నిర్వచిస్తాం.

\begin{defn}
$\mModel{M} =\tuple{W, R, V}$ అనుకుందాం;
$\Gamma$ \tetoken{ఉపసూత్రాల}{!!{formula}s} పరంగా సంవృతమై ఉండనివ్వండి.
$\Gamma$లోని ఒకే
\tetoken{సూత్రాలను}{!!{formula}s} సత్యం చేసే ఏ రెండు
లోకాలకైనా వర్తించేలా $W$పై $\equiv$ సంబంధాన్ని
నిర్వచిద్దాం. అంటే:
\[
u \equiv v \quad \text{అప్పుడు మరియు అప్పుడు మాత్రమే } \quad
\forall !A \in \Gamma : \mSat{M}{!A}[u] \Leftrightarrow \mSat{M}{!A}[v].
\]
లోకం~$w$ యొక్క తుల్యతా వర్గం~$[w]_\equiv$ను,
సంక్షిప్తంగా $[w]$ను, $w$కు $\equiv$-తుల్యమైన
అన్ని లోకాల సమితిగా నిర్వచిస్తాం:
\[
[w] = \Setabs{v}{v \equiv w}.
\]
\end{defn}

\begin{prop}
  $\mModel{M}$, $\Gamma$ ఇచ్చినప్పుడు, పైన
  నిర్వచించిన $\equiv$ తుల్యతా సంబంధం; అంటే అది
  స్వావర్తన, సౌష్ఠవ, సంక్రామక ధర్మాలు కలిగి ఉంటుంది.
\end{prop}

\begin{proof}
  $w$, $\Gamma$లోని ఏ
  \tetoken{సూత్రాలను}{!!{formula}s} తన వద్ద సత్యం
  చేస్తుందో అవే తన వద్ద సత్యం చేస్తుంది; కనుక
  $\equiv$ స్వావర్తనం. $u$, $\Gamma$లోని ఏ
  \tetoken{సూత్రాలను}{!!{formula}s} సత్యం చేస్తుందో
  $v$ అవే సత్యం చేస్తే, $v$ ఏవి సత్యం చేస్తుందో
  $u$ కూడా అవే సత్యం చేస్తుంది; కనుక సౌష్ఠవం.
  అలాగే $u$, $\Gamma$లోని ఏ
  \tetoken{సూత్రాలను}{!!{formula}s} సత్యం చేస్తుందో
  $v$ అవే చేస్తే, $v$ చేసే వాటినే $w$ కూడా చేస్తే,
  $u$ చేసే \tetoken{సూత్రాలనే}{!!{formula}s} $w$ కూడా సత్యం చేస్తుంది; కనుక
  సంక్రామకం.
\end{proof}

ఏ తుల్యతా సంబంధంలాగే, $\equiv$ కూడా~$W$ను
\emph{విభాగాలు}గా విభజిస్తుంది: అవి $W$కు చెందిన,
జతలవారీగా వియుక్తమైన ఉపసమితులు; అన్నీ కలిపితే
మొత్తం~$W$ను ఆవరిస్తాయి. ప్రతి $w \in W$ ఒక
విభాగానికి, అంటే~$[w]$కు,
\tetoken{మూలకం}{!!a{element}}; ఎందుకంటే $w \equiv w$.
అందువల్ల $[w]$ విభాగాలు మొత్తం~$W$ను ఆవరిస్తాయి.
అవి జతలవారీగా వియుక్తం: $u \in [w]$, $u \in [v]$
అయితే $u \equiv w$, $u \equiv v$; సౌష్ఠవం,
సంక్రామకం వల్ల $w \equiv v$; కాబట్టి $[w] = [v]$.

\end{document}
